\documentclass{openreader}
\title{Ἰωάννου γʹ}
\date{}
\setmainlanguage{english}
\setmainfont{Charis SIL}
\setotherlanguage{greek}
\newfontfamily\greekfont[Script=Greek]{SBL BibLit}
\begin{document}
\ORBselectlanguage{greek}
\maketitle
\raggedbottom 
\fontsize{16pt}{24pt}\selectfont

\ChapterHeader{1}{Chapter 1}

\MainTextVerseMark{1}{1}Ὁ πρεσβύτερος Γαΐῳ\FnParseFormGloss{a}{dative masculine singular}{Γάϊος}{Gaius} τῷ ἀγαπητῷ, ὃν ἐγὼ ἀγαπῶ\FnParse{b}{present active indicative 1st singular} ἐν ἀληθείᾳ. 
\MainTextVerseMark{1}{2}Ἀγαπητέ, περὶ πάντων εὔχομαί\FnParseFormGloss{a}{present middle indicative 1st singular}{εὔχομαι}{to pray for} σε εὐοδοῦσθαι\FnParseFormGloss{b}{present passive infinitive}{εὐοδόομαι}{to get along with} καὶ ὑγιαίνειν,\FnParseFormGloss{c}{present active infinitive}{ὑγιαίνω}{to be healthy} καθὼς εὐοδοῦταί\FnParseFormGloss{d}{present passive indicative 3rd singular}{εὐοδόομαι}{to get along with} σου ἡ ψυχή. 
\MainTextVerseMark{1}{3}ἐχάρην\FnParse{a}{aorist passive indicative 1st singular} γὰρ λίαν\FnFormGloss{b}{λίαν}{very much} ἐρχομένων\FnParse{c}{present middle participle genitive masculine plural} ἀδελφῶν καὶ μαρτυρούντων\FnParse{d}{present active participle genitive masculine plural} σου τῇ ἀληθείᾳ, καθὼς σὺ ἐν ἀληθείᾳ περιπατεῖς.\FnParse{e}{present active indicative 2nd singular} 
\MainTextVerseMark{1}{4}μειζοτέραν\FnParseFormGloss{a}{accusative feminine singular}{μειζότερος}{greater} τούτων οὐκ ἔχω\FnParse{b}{present active indicative 1st singular} ⸀χαράν, ἵνα ἀκούω\FnParse{c}{present active subjunctive 1st singular} τὰ ἐμὰ τέκνα ἐν ⸀τῇ ἀληθείᾳ περιπατοῦντα.\FnParse{d}{present active participle accusative neuter plural} 
\MainTextVerseMark{1}{5}Ἀγαπητέ, πιστὸν ποιεῖς\FnParse{a}{present active indicative 2nd singular} ὃ ἐὰν ἐργάσῃ\FnParse{b}{aorist middle subjunctive 2nd singular} εἰς τοὺς ἀδελφοὺς καὶ ⸀τοῦτο ξένους,\FnParseFormGloss{c}{accusative masculine plural}{ξένος}{strange} 
\MainTextVerseMark{1}{6}οἳ ἐμαρτύρησάν\FnParse{a}{aorist active indicative 3rd plural} σου τῇ ἀγάπῃ ἐνώπιον ἐκκλησίας, οὓς καλῶς ποιήσεις\FnParse{b}{future active indicative 2nd singular} προπέμψας\FnParseFormGloss{c}{aorist active participle nominative masculine singular}{προπέμπω}{to accompany} ἀξίως\FnFormGloss{d}{ἀξίως}{worthily} τοῦ θεοῦ· 
\MainTextVerseMark{1}{7}ὑπὲρ γὰρ τοῦ ὀνόματος ἐξῆλθον\FnParse{a}{aorist active indicative 3rd plural} μηδὲν λαμβάνοντες\FnParse{b}{present active participle nominative masculine plural} ἀπὸ τῶν ⸀ἐθνικῶν.\FnParseFormGloss{c}{genitive masculine plural}{ἐθνικός}{pagan} 
\MainTextVerseMark{1}{8}ἡμεῖς οὖν ὀφείλομεν\FnParse{a}{present active indicative 1st plural} ⸀ὑπολαμβάνειν\FnParseFormGloss{b}{present active infinitive}{ὑπολαμβάνω}{to take up} τοὺς τοιούτους, ἵνα συνεργοὶ\FnParseFormGloss{c}{nominative masculine plural}{συνεργός}{fellow worker} γινώμεθα\FnParse{d}{present middle subjunctive 1st plural} τῇ ἀληθείᾳ. 
\MainTextVerseMark{1}{9}Ἔγραψά\FnParse{a}{aorist active indicative 1st singular} ⸀τι τῇ ἐκκλησίᾳ· ἀλλ’ ὁ φιλοπρωτεύων\FnParseFormGloss{b}{present active participle nominative masculine singular}{φιλοπρωτεύω}{to love to be first} αὐτῶν Διοτρέφης\FnParseFormGloss{c}{nominative masculine singular}{Διοτρέφης}{Diotrephes} οὐκ ἐπιδέχεται\FnParseFormGloss{d}{present middle indicative 3rd singular}{ἐπιδέχομαι}{to welcome} ἡμᾶς. 
\MainTextVerseMark{1}{10}διὰ τοῦτο, ἐὰν ἔλθω,\FnParse{a}{aorist active subjunctive 1st singular} ὑπομνήσω\FnParseFormGloss{b}{future active indicative 1st singular}{ὑπομιμνῄσκω}{to remind} αὐτοῦ τὰ ἔργα ἃ ποιεῖ,\FnParse{c}{present active indicative 3rd singular} λόγοις πονηροῖς φλυαρῶν\FnParseFormGloss{d}{present active participle nominative masculine singular}{φλυαρέω}{to gossip} ἡμᾶς, καὶ μὴ ἀρκούμενος\FnParseFormGloss{e}{present passive participle nominative masculine singular}{ἀρκέω}{to be content} ἐπὶ τούτοις οὔτε αὐτὸς ἐπιδέχεται\FnParseFormGloss{f}{present middle indicative 3rd singular}{ἐπιδέχομαι}{to welcome} τοὺς ἀδελφοὺς καὶ τοὺς βουλομένους\FnParse{g}{present middle participle accusative masculine plural} κωλύει\FnParseFormGloss{h}{present active indicative 3rd singular}{κωλύω}{to hinder} καὶ ἐκ τῆς ἐκκλησίας ἐκβάλλει.\FnParse{i}{present active indicative 3rd singular} 
\MainTextVerseMark{1}{11}Ἀγαπητέ, μὴ μιμοῦ\FnParseFormGloss{a}{present middle imperative 2nd singular}{μιμέομαι}{to imitate} τὸ κακὸν ἀλλὰ τὸ ἀγαθόν. ὁ ἀγαθοποιῶν\FnParseFormGloss{b}{present active participle nominative masculine singular}{ἀγαθοποιέω}{to do good} ἐκ τοῦ θεοῦ ἐστιν·\FnParse{c}{present active indicative 3rd singular} ὁ κακοποιῶν\FnParseFormGloss{d}{present active participle nominative masculine singular}{κακοποιέω}{to do evil} οὐχ ἑώρακεν\FnParse{e}{perfect active indicative 3rd singular} τὸν θεόν. 
\MainTextVerseMark{1}{12}Δημητρίῳ\FnParseFormGloss{a}{dative masculine singular}{Δημήτριος}{Demetrius} μεμαρτύρηται\FnParse{b}{perfect passive indicative 3rd singular} ὑπὸ πάντων καὶ ὑπὸ αὐτῆς τῆς ἀληθείας· καὶ ἡμεῖς δὲ μαρτυροῦμεν,\FnParse{c}{present active indicative 1st plural} καὶ ⸀οἶδας\FnParse{d}{perfect active indicative 2nd singular} ὅτι ἡ μαρτυρία ἡμῶν ἀληθής\FnParseFormGloss{e}{nominative feminine singular}{ἀληθής}{true} ἐστιν.\FnParse{f}{present active indicative 3rd singular} 
\MainTextVerseMark{1}{13}Πολλὰ εἶχον\FnParse{a}{imperfect active indicative 1st singular} ⸂γράψαι\FnParse{b}{aorist active infinitive} σοι⸃, ἀλλ’ οὐ θέλω\FnParse{c}{present active indicative 1st singular} διὰ μέλανος\FnParseFormGloss{d}{genitive neuter singular}{μέλαν}{the color black} καὶ καλάμου\FnParseFormGloss{e}{genitive masculine singular}{κάλαμος}{reed} σοι ⸀γράφειν·\FnParse{f}{present active infinitive} 
\MainTextVerseMark{1}{14}ἐλπίζω\FnParse{a}{present active indicative 1st singular} δὲ εὐθέως ⸂σε ἰδεῖν⸃,\FnParse{b}{aorist active infinitive} καὶ στόμα πρὸς στόμα λαλήσομεν.\FnParse{c}{future active indicative 1st plural} 
\MainTextVerseMark{1}{15}Εἰρήνη σοι. ἀσπάζονταί\FnParse{a}{present middle indicative 3rd plural} σε οἱ φίλοι.\FnParseFormGloss{b}{nominative masculine plural}{φίλος}{friendly} ἀσπάζου\FnParse{c}{present middle imperative 2nd singular} τοὺς φίλους\FnParseFormGloss{d}{accusative masculine plural}{φίλος}{friendly} κατ’ ὄνομα. 
\end{document}